\chapter*{Résumé}

À mesure que les systèmes photovoltaïques (PV) deviennent un élément central des systèmes énergétiques modernes, une fiabilité durable dans des conditions réelles d'exploitation est essentielle à leur valeur économique. Les centrales PV sont exposées à des défauts qui peuvent être invisibles à l'inspection directe et se manifestent principalement par des écarts dans les mesures électriques et météorologiques. Une détection et un diagnostic précoces sont donc nécessaires pour préserver les performances, limiter la dégradation et permettre une supervision à grande échelle.

Cette thèse développe un système frugal de détection et de diagnostic des défauts pour les centrales PV à partir de mesures réelles de terrain, adapté à un déploiement embarqué en périphérie. Le système suit une structure à deux étapes : la détection d'anomalies identifie les écarts par rapport au comportement normal, tandis que la classification des défauts détermine le type de défaut correspondant. Une interface de surveillance en temps réel facilite le suivi opérationnel.

La principale contribution méthodologique est PC-Flow, une architecture de flux normalisant conditionné par la physique, conçue pour une détection d'anomalies précise et hautement efficace. Avec moins de 25\,000 paramètres, PC-Flow atteint un PR-AUC binaire de 0{,}972 et réalise l'inférence en 970~$\mu$s au 99\textsuperscript{e} centile. Pour la classification des défauts, le classifieur ExtraTrees proposé atteint un macro F1-score de 0{,}959 et une précision équilibrée de 94{,}3\,\%. Sur la plateforme embarquée, le pipeline complet reste inférieur à 31~ms au 99\textsuperscript{e} centile.

Ces résultats établissent qu'une supervision des défauts précise et à faible latence est réalisable directement au niveau de la centrale, réduisant le délai entre l'apparition d'un défaut et la réponse de l'opérateur, limitant les pertes de production cumulées et rendant la surveillance PV continue viable dans les conditions de terrain où elle est la plus nécessaire.

\vspace{1cm}
\noindent\textbf{Mots-clés :} systèmes photovoltaïques, détection et diagnostic de défauts, flux normalisants, modèles informés par la physique, inférence embarquée.
